\section{The extension, not another architecture}
\label{sec:delta}

\subsection{What the repository already supplies}
The baseline for this article is the inspected Forge revision
\code{58ea206bd7ad} (full identifier in Appendix~\ref{app:status}), not the
initial proposal to improve \code{grind}. Its README and provenance ledger describe a merge of
27 proposals across three rounds. The merged design already includes structural
proof planning, generalization, synthesized witnesses and invariants,
observable-space and ideal closure, finite-algebra covers, tree summaries and
observable quotients, summation and Ore transport, noncommutative algebra,
analytic enclosures, quantitative probabilistic workers, games, and unbounded
pushdown summaries~\cite{forge}. A generic instruction to add induction, a CAS,
an e-graph, or a certificate boundary would therefore be a repetition.

Nor should a new report claim that none of Forge's Lean has ever compiled. The
README records 18 core-only files elaborated by the merge, even though the
contributing proposal environments did not run Lean. Those elaborations do not
establish an implemented \code{forge} tactic or a verified family of certificate
checkers. The present environment likewise has no Lean executable; that is a
limitation of this delivery, not a revision of the repository's own evidence.

The audit used the merged article's closure section and provenance inventory,
the README, and the concrete runtime, logical-contract, and finite-cover Lean
files. The Leant/Djex comparison used their current revision information and
indexing-composition diagnostics. It was not a line-by-line review of every
archived proposal. In particular, the empty results returned by the repository
code-search endpoint were not treated as evidence of absence: even a known
baseline term returned no match. The delta claims below are relative to the
inspected merged design, with the nearest existing capability named explicitly.

\begin{longtable}{@{}>{\raggedright\arraybackslash}p{.22\textwidth}>{\raggedright\arraybackslash}p{.32\textwidth}>{\raggedright\arraybackslash}p{.40\textwidth}@{}}
\toprule
Addition & Nearest existing Forge mechanism & Actual additional work \\
\midrule\endhead
Counter antichains & Numeric invariants, ideal closure, and finite covers &
Compute exact backward coverability in $Q\times\N^d$; retain consuming guards;
certify all members of nonnegative affine initial families. \\
Equality orbits & Finite-algebra covers and observable quotients &
Construct the finite quotient for equality-only data over infinitely many atoms;
prove joint support and fresh-case coverage; recover concrete witnesses. \\
Orbit--counter product & Separate infinite-state workers &
Compile finite name orbits to controls while retaining unbounded counters, and
check directly against the original mixed-rule semantics. \\
Candidate discharge & Existing typed synthesis and behavioral filtering &
Use these exact semantic summaries after candidate reconstruction, with a
proved source bridge, rather than promote finite observations to a theorem. \\
\bottomrule
\end{longtable}

\subsection{What makes the new workers useful}
A finite set of \emph{concrete states} cannot cover an arbitrarily large pool of
waiting clients. A finite set of \emph{minimal bad configurations} can describe
all markings from which a violation is possible. Similarly, finitely many
concrete keys do not represent every name a program might receive, but finitely
many equality patterns relative to its live registers do.

These are different reasons for finiteness. For counters, finiteness comes from
a well quasi-order on natural vectors. For names, it comes from quotienting by
renamings that preserve everything the program can observe. Combining them
requires a precise intermediate language; it is not justified merely by saying
that both components separately have decidable fragments.

The intended Forge results are ordinary mathematical theorems, for example
\[
 \forall n\in\N,\ \forall\text{ finite executions from }s_0(n),
 \qquad \mathrm{active}\le1,
\]
or
\[
 \forall w\in(\mathrm{Tag}\times\Atoms)^*,
 \qquad \mathrm{outputs}_A(w)=\mathrm{outputs}_B(w).
\]
The quantifiers range over all initial resource sizes and all finite input
lengths. Neither is replaced by a test bound. The finite object returned by a
worker explains why those quantifiers can be discharged.

\subsection{Historical versus project novelty}
The mathematical foundations are established: monotone systems and
well-quasi-order methods for coverability~\cite{abdulla,finkel}, and automata
over nominal sets for data alphabets with symmetries~\cite{nominal}. We do not
rename these results and claim to have invented them. The contribution here is
a specific fragment choice, exact certificate formats, an initial-family
interface, an independently computed replay path, a mixed-rule frontend, and
an integration plan tied to Forge's actual types.

This distinction also prevents an exaggerated expressiveness claim. A hand-made
induction invariant may prove the mutual-exclusion example, and the repository's
polynomial machinery could express some of the same properties. A complete
specialized worker is still useful when it discovers the finite invariant
without choosing a template, emits a counterexample when the property is false,
and explains exactly when its answer is exhaustive. No delivered experiment
compares it with existing Lean tactics.

\subsection{Three deliberately narrow interfaces}
The executable workers accept mathematical intermediate representations, not
arbitrary Python or Lean programs. Their admission conditions are part of their
meaning. A counter system has nonnegative integer consumption and production,
upward-closed bad sets, and no hidden tests. An equality machine has only finite
control, fixed constants, equality tests, and copying registers. A mixed system
has both components, but counter updates do not depend on the numerical identity
of an atom and do not allocate an unbounded map of counters indexed by names.

The source program remains outside this small language until a proved bridge
connects its execution to the representation. That bridge is where most Lean
implementation effort belongs. Another search heuristic cannot make an
unjustified translation sound.
